% ============ 第1章 §1.10 高维情形的推广 (原书 pp.31-32) ============

\section{高维情形的推广}\label{sec:1.10}
\index{0symbolvecX@$\vec{X}_{k}$}

\index{0symbolDkm@$D_{k}^{m}$}

\index{0symbolE@$E$}

\index{0symbolEk2@$E_{k}$ (高维)}

\index{0symbolEkm@$E_{k}^{m}$}












用整数分量向量作指标编号的数列称为多重数列. 这里取 N, n 为自然数, 考虑由 $2^{nN}$ 个实数组成的多重数列 $a = \{a_{\vec{j}}\}_{\vec{j}\in I(N,0)^{n}}$. 由于这样的数列由 n 维向量组成, 称为 n 重数列或 n 维数列. 区间 $I(N,0)$ 是在式~\eqref{eq:1.14} 中令 j = N, k = 0 定义的. 暂且固定 $k = 0, 1, \ldots, N$. 迄今以 n = 1, 除 $D_{N}$ 外考虑了 $E_{k}, D_{k}$ 等算子, 将把它们作用于第 $m\,(=1, 2, \ldots, n)$ 分量所得的算子记为 $E_{k}^{m}, D_{k}^{m}$. 也就是说, 例如, 对 $\vec{j} = (j_{1}, j_{2}, \ldots, j_{n})$, 定义 $E_{k}^{m}(a) = \{E_{k}^{m}(a)_{\vec{j}}\}_{\vec{j}\in I(N,0)^{n}}$ 为
\[
E_{k}^{m}(a)_{\vec{j}} = \frac{1}{2^{k}}\sum_{l = 2^{k}[2^{-k}j_{m}]}^{2^{k}[2^{-k}j_{m}] + 2^{k} - 1}a_{(j_{1}, j_{2}, \ldots, j_{m-1}, l, j_{m+1}, \ldots, j_{n})}.
\]

继续设 $k = 0, 1, \ldots, N$. 设给定由 D, E 中某一字母组成的 n 个字母的序列 $\vec{X} = (X^{1}, X^{2}, \ldots, X^{n})$. 在这样的序列中, 仅 E 连续出现 n 次的序列记为 E. 一般地, 当 $X^{j} = D$ 时令 $Y^{j} = D_{k}^{j}$, 当 $X^{j} = E$ 时令 $Y^{j} = E_{k+1}^{j}$, 令由 $\vec{X}_{k} = Y^{1} \circ Y^{2} \circ \cdots \circ Y^{n}$ 得到的算子为 $\vec{X}_{k}$. 例如, 当 n = 3 时,
\[
(D, D, E)_{k} = D_{k}^{1} \circ D_{k}^{2} \circ E_{k+1}^{3}, \qquad (E, D, E)_{k} = E_{k+1}^{1} \circ D_{k}^{2} \circ E_{k+1}^{3}.
\]
又, $\mathbf{E}_{k} = (E, E, \ldots, E)_{k}$. 一般地, 对多重数列 a, 由于 $\mathbf{E}_{0}(a) = a$, 分解为
\[
a = \mathbf{E}_{N}(a) + \sum_{k=0}^{N-1}(\mathbf{E}_{k}(a) - \mathbf{E}_{k+1}(a)).
\]
进而, 由于 $D_{k} = E_{k} - E_{k+1}$, 可分解为
\[
\mathbf{E}_{k} - \mathbf{E}_{k+1} = \sum_{\vec{X}\in\{(D,E)\}^{n}\setminus\{\mathbf{E}\}}\vec{X}_{k}.
\]
基于这一点, 对 $\vec{j} \in I(N,0)^{n}$, 定义实数 $g(a)_{\vec{j}} \geq 0$ 为
\[
g(a)_{\vec{j}} = \left(|\mathbf{E}_{N}(a)_{\vec{j}}|^{2} + \sum_{k=1}^{N}\sum_{\vec{X}\in\{(D,E)\}^{n}\setminus\{\mathbf{E}\}}|\vec{X}_{k-1}(a)_{\vec{j}}|^{2}\right)^{1/2}.
\]
记为 $g(a) = \{g(a)_{\vec{j}}\}_{\vec{j}\in I(N,0)^{n}}$. 用与 n = 1 情形相同的证明方法, 可以证明以下结论.

\begin{theorem}\label{thm:1.14}
设 $n, N \in \mathbb{N}$. 存在仅依赖于 p, n 的常数 $C_{p,n}$, 当 $1 < p < \infty$ 时, 对由 $2^{nN}$ 个分量组成的多重数列 $a = \{a_{\vec{j}}\}_{\vec{j}\in I(N,0)^{n}}$, 成立
\[
C_{p,n}^{-1}\|a\|_{\ell^{p}} \leq \|g(a)\|_{\ell^{p}} \leq C_{p,n}\|a\|_{\ell^{p}}.
\]
\end{theorem}

\begin{exercise}\label{exer:1.27}
设 $k, k' = 0, 1, \ldots, N$, n = 2, 证明 $E_{k}^{1} \circ E_{k'}^{2} = E_{k'}^{2} \circ E_{k}^{1}$.
\end{exercise}
